\documentclass[11pt]{article}
\title{Sovereign AI Federation v2}
\date{}
\author{SIGMA AI-assisted research synthesis}
\usepackage[T1]{fontenc}
\usepackage[utf8]{inputenc}
\usepackage[margin=1in]{geometry}
\usepackage{amsmath,amssymb,array,booktabs,longtable,pdflscape}
\usepackage{microtype}
\usepackage{titlesec}
\emergencystretch=8em
\raggedbottom
\linespread{1.07}
\usepackage{xcolor}
\definecolor{blue}{HTML}{24485E}
\titleformat{\section}{\Large\bfseries\sffamily\color{blue}}{\thesection}{0.8em}{}
\titleformat{\subsection}{\large\bfseries\sffamily}{\thesubsection}{0.7em}{}
\setlength{\parskip}{.6em}
\setlength{\parindent}{0pt}
\begin{document}
\maketitle
\begin{center}\fbox{\parbox{.92\linewidth}{\textbf{Needs review: TeX verification is incomplete.} Review the archived proofreading diagnostics before relying on the typesetting.}}\end{center}
\begin{center}\small\textit{AI-generated research synthesis \textperiodcentered{} SIGMA}\end{center}

\begin{abstract}
The investigation asks: Can universities with uneven GPU capacity provide useful, reliable AI research access while retaining control of sensitive data and staying within hard budgets? The saved research reports a provisional finding, not independently verified here: Structured argument steps are recorded; review their evidence and objections. The analysis distinguishes recorded arguments from demonstrated results, identifies remaining gaps, and does not claim experiments or proof steps absent from the frozen record.
\end{abstract}
\section*{About this paper}
This paper was generated by SIGMA, an AI-assisted research system, from the project's recorded investigations, findings, and supporting evidence.
The analysis and conclusions reflect an automated synthesis of that research. Although SIGMA incorporates evidence tracking, critical review, and consistency checks, the resulting paper has not necessarily undergone independent verification or peer review.
Readers should evaluate its arguments, supporting evidence, and stated limitations before relying on its conclusions.
Research style: Distinguish algorithmic guarantees from tests, benchmarks and conjectures. Document counterexamples, complexity assumptions and reproducibility.
\section*{Evidence coverage limitation}
The evidence index omitted 9706 material saved research field fragment(s) across 10 source(s). These fragments were not reviewed for this paper. Interpretations that could depend on omitted material remain provisional; the complete frozen source snapshot and omission ledger are retained with the generation record.
\textit{Claims discussed: }C1, C2
\section{Introduction and problem statement}
The investigation asks: Can universities with uneven GPU capacity provide useful, reliable AI research access while retaining control of sensitive data and staying within hard budgets? Compare four baselines: (1) campus-local inference only, (2) approved frontier API only, (3) a campus gateway with separately approved external access, and (4) a proposed cooperative federation among campuses. Treat federation and frontier routing as hypothetical extensions, not existing MindRouter features.

Ground the investigation in MindRouter's published campus gateway: protocol translation across OpenAI/Ollama/Anthropic clients and Ollama/vLLM backends; WDRR fair-share scheduling; per-user quotas; audit logging; and GPU telemetry. Its documentation describes an on-campus service, not a multi-institution federation. Start with its documentation (https://mindrouter.ai/document

\textit{Claims discussed: }C1
\section{Methods and evidence}
This synthesis organizes the project's frozen research records and their saved evidence references. It does not claim new experiments.

\section{Background and research context}
The research question is whether shared AI infrastructure can expand useful university research access without transferring control over sensitive data or weakening hard budget constraints. The comparison concerns four alternatives: campus-local inference, approved frontier APIs, a campus gateway with separately approved external access, and a proposed cooperative federation. Federation is therefore an option to evaluate against a controlled campus service, not an assumed solution to uneven GPU capacity. Its relevant contribution would be additional policy-permitted research capability, rather than resource sharing alone. [C1]

MindRouter provides the starting campus-gateway context. The supplied project description attributes protocol translation across OpenAI, Ollama and Anthropic clients and Ollama/vLLM backends, WDRR fair-share scheduling, per-user quotas, audit logging and GPU telemetry to its published on-campus service. The starting references are the documentation at https://mindrouter.ai/documentation.html and the PEARC\textquoteright{}26 short paper at https://ai4ra.uidaho.edu/wp-content/uploads/2026/04/Sheneman\_MindRouter\_PEARC26\_Short\_Paper.pdf.

However, neither reference was retrieved in the supplied research record: 2026-10-31 is a project-reported historical review date, not an independently established access date, and the feature descriptions remain attributed claims rather than inspected or measured capabilities. Cross-campus federation, frontier routing and the proposed financial enforcement controls are consequently treated as hypothetical extensions, not existing MindRouter features. [C1]

The alternatives differ in execution domains and institutional dependencies, as specified in the comparison table. In particular, separately approved external access need not imply automatic gateway routing: an explicitly authorized portal or account workflow is a distinct baseline. Adding federation to that baseline introduces partner-campus execution without, by itself, enlarging the set of destinations authorized for any particular task. This distinction makes it possible to ask whether cooperation adds a capability unavailable through local service and separately approved external access, rather than attributing every external-model benefit to federation. [C1]

The study defines sovereignty through retained campus authority over identity, classification, approved processing domains, retention, audit requirements, contracts and quotas. Its proposed policy boundary covers derived prompts, outputs, logs, embeddings and caches as well as original inputs; restricted, student, patient, controlled, sponsor-restricted and unpublished content remains within approved domains by default. Ambiguous classification supplies no execution destination until resolved. Quality estimates, available capacity and incremental cost can rank eligible destinations, but overload, outages, retries, timeouts or exhausted budgets cannot create permission to use an otherwise prohibited destination. [C1]

Useful access is defined at the level of an authorized root research task, not an individual model call. Classroom coding, routine analysis, demanding scientific reasoning, multimodal or long-context work, agentic workflows and reproducible batches provide the proposed workload scope, while sensitivity is a classification that must also be tracked within task categories. A useful completion must satisfy prespecified quality, deadline and provenance requirements; tool calls and retries remain descendants of the same task, and a correct policy refusal is recorded separately from delivered research.

This accounting basis prevents multi-call workflows from inflating completion counts and keeps restricted-demand coverage visible rather than allowing aggregate performance to stand in for every classification stratum. [C1]

The proposed comparison separates permission to execute from the technical opportunity to complete a task. Relevant technical conditions include GPU memory and model/context fit, compatible capacity, queueing, cold starts, maintenance, peak demand and the allocation of service between interactive and batch work. Latency, throughput, availability, fairness and starvation therefore belong alongside task quality and provenance/version reproducibility in the evaluation scope; scheduling weights alone are not treated as service guarantees.

Cost comparisons likewise require a common workload and horizon, with sunk, incremental and fully loaded accounts kept distinct: local ownership and operations include capital, replacement, idle capacity, energy, cooling, storage, networking, support, procurement and staff, whereas paid workflows include their constituent calls, retries, tools, embeddings, egress and contractual commitments. [C1]

The conditional design also makes its authority and accounting assumptions explicit. It assumes a durable, linearizable home-campus authority, trusted binding of classifications to immutable requests, enforceable nested liabilities and complete mediation of invocation and spending paths, including separately approved accounts. Its decision-time authorization semantics permit an earlier committed grant to transmit after a later revocation; they do not establish a stronger prohibition on subsequent delivery or execution, and uncertain transmission is not repaired by reissuing a consumed grant. These assumptions define the proposed extension\textquoteright{}s scope rather than demonstrated deployment properties.

The evidence thus frames the central question more precisely: federation must add usable, authorized research capability over the gateway-plus-approved-access baseline while preserving the same campus authority and accounting obligations; nominal pooled capacity is not sufficient to establish that advantage. [C1]

\begin{longtable}{@{}>{\raggedright\arraybackslash}p{\dimexpr(\linewidth-2\tabcolsep*2)/126*42\relax}>{\raggedright\arraybackslash}p{\dimexpr(\linewidth-2\tabcolsep*2)/126*42\relax}>{\raggedright\arraybackslash}p{\dimexpr(\linewidth-2\tabcolsep*2)/126*42\relax}@{}}
\caption{Comparison scope for the four alternatives. These are architectural definitions and dependencies, not measured rankings or verified federation capabilities. [C1]}\\
\toprule
\textbf{Alternative} & \textbf{Execution scope} & \textbf{Dependencies relevant to the comparison}\\
\midrule
\endfirsthead
\toprule
\textbf{Alternative} & \textbf{Execution scope} & \textbf{Dependencies relevant to the comparison}\\
\midrule
\endhead
\midrule
\multicolumn{3}{r}{\footnotesize Continued on next page}\\
\endfoot
\bottomrule
\endlastfoot
Campus-local inference only & Campus-hosted models and infrastructure & Local model fit, capacity, scheduling, maintenance and campus operations\\
Approved frontier API only & Explicitly approved provider destinations & Data eligibility, provider availability and rejection, paid caps and model-version availability\\
Campus gateway plus separately approved external access & Campus inference and a distinct authorized external workflow & Local identity, audit and quotas; separate external approval; complete financial mediation if a hard-cap guarantee is invoked\\
Proposed cooperative federation & Gateway-plus-approved-access baseline extended with approved partner-campus execution & Compatible partner capacity, identity mapping, WAN connectivity, recipient restrictions, incident coordination and cross-campus agreements\\
\end{longtable}
\par\smallskip{\footnotesize Evidence: [Research session 1, artifact 3], [Research session 1, artifact 2]; claims: C1}
\textit{Claims discussed: }C1
\section{Recorded findings and observations}
The recorded analysis makes authorized, useful root-task completion\textemdash{}not inference calls, tokens, or nominal GPU capacity\textemdash{}the unit for comparing research access. Its substantive contribution to claim C2 is to separate policy eligibility, actual admission, and successful delivery, and to show that aggregate success can conceal failure for campus-only sensitive work. The supporting results are conditional deductions and hypothetical accounting examples: the supplied evidence contains no measured four-baseline comparison, and neither the starting MindRouter sources nor current provider prices and institutional permissions were independently inspected.

Federation and frontier routing therefore remain proposed extensions rather than established MindRouter capabilities. [C2; [Research session 1, artifact 2]; [Research session 1, artifact 8]]

A root task is one research objective with a fixed input snapshot, requested artifact, quality rubric, deadline, classification, provenance requirements, and resource envelope. Agent calls, tools, retries, critique exchanges, and reproducibility repeats are descendants of that objective, not additional completed tasks. In the recorded acceptance rule, a task succeeds only when its final artifact satisfies the quality rubric, required provenance, authorized processing, and end-to-end deadline jointly.

A correct policy refusal is recorded as safe behavior but not as delivered research; refusals and policy-ineligible demand consequently need separate reporting. [C2; [Research session 1, artifact 6]; [Research session 1, artifact 2]]

This definition requires task-specific acceptance rather than a common model score. The recorded proposals use held-out tests for coding, known-result and numerical checks for analysis, expert assessment for scientific reasoning, labeled references for multimodal tasks, and root-goal plus tool-boundary checks for agentic workflows. Reproducible batches additionally require a registered rerun or numerical-tolerance criterion; a model-version label alone does not establish reproducibility. Delivery latency begins at submission and includes approval handling, queues, cold starts, tools, retries, and recovery, rather than ending at the first streamed token. [C2; [Research session 1, artifact 6]]

The numerical deadlines are not consistent across the recorded proposals and cannot be treated as a settled service specification. One proposal assigns coding a deadline of $30$ seconds and agentic workflows $900$ seconds; another assigns coding $120$ seconds and agentic workflows $60$ minutes. Both describe illustrative targets awaiting approval, not observations. Because deadline compliance is part of usefulness, changing these values changes the outcome being compared even when the generated artifact is unchanged. [C2; [Research session 1, artifact 2]; [Research session 1, artifact 6]]

The stratum-specific formulation makes the comparison precise. Let $N_s$ denote offered root tasks in classification-and-task stratum $s$, and let $n_{bds}$ denote distinct roots assigned to destination $d$ under baseline $b$, with $\sum_d n_{bds}\leq N_s$ and unauthorized assignments constrained to $n_{bds}=0$. If $q_{bds}$ is the average marginal probability that those admissions jointly satisfy authorization, quality, deadline, and provenance requirements, expected useful completions are $\mu_b=\sum_{s,d}n_{bds}q_{bds}$.

This expectation identity requires no independence assumption, but its probabilities concern the selected workload, load, and configuration rather than universal model success rates. [C2; [Research session 1, artifact 8]]

The most consequential counterexample concerns sensitive-stratum coverage. For $200$ sensitive and $800$ other local admissions, aggregate expected usefulness of $0.80$ establishes only $200q_S+800q_N=800$. The assignment $q_S=0$ and $q_N=1$ satisfies that equation while delivering no expected sensitive completions; the same counterexample can be expressed as exact scenario counts. Thus neither aggregate usefulness nor sensitive-first admission establishes useful sensitive service, and the recorded analysis explicitly withdraws the earlier unqualified sensitive-feasibility inference. [C2; [Research session 1, artifact 8]]

The repaired example adds the missing premise rather than inferring it from the aggregate: both local strata have hypothetical joint usefulness of $0.80$, while selected API admissions have hypothetical usefulness of $0.90$. Its common monthly cohort contains $2{,}000$ roots, including $200$ sensitive campus-only roots, and assumes deadline-feasible local admission of $1{,}000$ roots. The illustrative acceptance gates require $160$ expected useful sensitive roots and $1{,}400$ expected useful roots overall.

These assumptions produce the comparison in the table, whose federation row is deliberately restricted to the recorded case of no additional usable peer service and no offsetting benefit. [C2; [Research session 1, artifact 8]]

The repaired scenario exposes different failure modes among the alternatives. Campus-local service misses the overall expected-count gate because admission capacity is limited, whereas API-only service cannot satisfy the campus-only sensitive gate regardless of its aggregate completion count. For the gateway with separately approved external access, let $B$ be the API allocation and $r>0$ an enforceable maximum liability per admitted API root; then $n_G=\min(1000,\lfloor B/r\rfloor)$ and $\mu_G=800+0.90n_G$. The two expected-count gates hold precisely when $\lfloor B/r\rfloor\geq667$, given the stated local-stratum assumptions.

This threshold is an expectation calculation, not a guarantee that the realized workload meets either gate. [C2; [Research session 1, artifact 8]]

The reservation-based admission count also has a narrower interpretation than an API price calculation. It assumes complete financial mediation, durable atomic reservations across applicable quotas, nonbinding subsidiary quotas, no opening liabilities or commitments, adequate API capacity, and no recycling of reservation savings within the fixed plan. The bound $r$ must cover workflow descendants and ancillary liabilities, including retries, tools, embeddings, and cache misses charged at uncached rates. An average observed charge or an optimistic cached-token estimate would not supply the enforceable maximum required by this argument. [C2; [Research session 1, artifact 8]]

Comparable denominators are equally important when moving from the scenario to a paired federation comparison. The recorded design distinguishes all requested tasks, authorized tasks, tasks permitted to undergo both workflows, and actual admissions in each workflow. Its primary comparison uses the same fixed paired population for the gateway and federation arms: model-fit failures, queue rejection, exhausted paid allocations, outages, and non-delivery remain outcomes rather than post hoc exclusions.

Tasks prohibited from paired evaluation are reported separately with reasons and arm-specific coverage, not exported for comparability; the paired population's share of all requested work exposes the resulting selection. [C2; [Research session 1, artifact 6]]

Within that fixed population, the paired effect is $\Delta=P(Y_F=1,Y_G=0)-P(Y_F=0,Y_G=1)$, where each outcome indicator applies the joint useful-completion rule. This measures federation gains net of tasks it harms, rather than counting only newly successful remote executions. The recorded proposal strengthens the benefit gate from a point estimate of at least $0.10$ with an interval excluding zero to a lower confidence bound exceeding $0.10$; evidence of any positive effect is not evidence that the effect exceeds that threshold.

The associated independent-pair calculations are analytical study-design results, with exact power and affordable enrollment unresolved, rather than observed federation benefits. [C2; [Research session 1, artifact 6]]

Finally, useful completions supply the denominator for the recorded economic result. For the same workload, horizon, acceptance rules, and cost basis, let the gateway cost and useful completions be $C_G>0$ and $S_G>0$, and let federation add net cost $\Delta C>0$ and net useful completions $\Delta S$, with $S_G+\Delta S>0$. Federation lowers cost per useful completion exactly when $\Delta S/S_G>\Delta C/C_G$; additional capacity or throughput alone is insufficient. In the restricted case of identical usable service opportunities, no offsetting federation benefit, and positive overhead, the gateway is cheaper.

The recorded findings therefore change the central comparison from whether federation pools more resources to whether it delivers enough additional authorized, timely, provenance-complete research outcomes in the required strata to justify its added cost. [C2; [Research session 1, artifact 2]; [Research session 1, artifact 8]]

\begin{landscape}
\small
\begin{longtable}{@{}>{\raggedright\arraybackslash}p{\dimexpr(\linewidth-2\tabcolsep*4)/189*42\relax}>{\raggedright\arraybackslash}p{\dimexpr(\linewidth-2\tabcolsep*4)/189*42\relax}>{\raggedright\arraybackslash}p{\dimexpr(\linewidth-2\tabcolsep*4)/189*31\relax}>{\raggedright\arraybackslash}p{\dimexpr(\linewidth-2\tabcolsep*4)/189*32\relax}>{\raggedright\arraybackslash}p{\dimexpr(\linewidth-2\tabcolsep*4)/189*42\relax}@{}}
\caption{Conditional monthly comparison for the repaired hypothetical cohort, not measured performance. Local admissions comprise $200$ sensitive and $800$ other roots, each with assumed joint usefulness $0.80$; selected API admissions have assumed joint usefulness $0.90$. With API allocation $B$ and enforceable per-root maximum liability $r>0$, define $n_A=\min(1800,\lfloor B/r\rfloor)$ and $n_G=\min(1000,\lfloor B/r\rfloor)$. The federation row applies only to the specified no-additional-usable-peer-service case.}\\
\toprule
\textbf{Baseline} & \textbf{Assumed admissions} & \textbf{Expected useful sensitive roots} & \textbf{Expected useful roots overall} & \textbf{Implication for illustrative expected-count gates}\\
\midrule
\endfirsthead
\toprule
\textbf{Baseline} & \textbf{Assumed admissions} & \textbf{Expected useful sensitive roots} & \textbf{Expected useful roots overall} & \textbf{Implication for illustrative expected-count gates}\\
\midrule
\endhead
\midrule
\multicolumn{5}{r}{\footnotesize Continued on next page}\\
\endfoot
\bottomrule
\endlastfoot
Campus-local only, $L$ & $1{,}000$ local roots & $160$ & $800$ & Meets sensitive gate; misses overall gate.\\
Approved API only, $A$ & $n_A$ API-eligible roots & $0$ & $0.90n_A$, at most $1{,}620$ & Cannot meet campus-only sensitive gate.\\
Campus gateway plus separately approved external access, $G$ & $1{,}000$ local roots plus $n_G$ distinct residual API roots & $160$ & $800+0.90n_G$, at most $1{,}700$ & Meets both gates exactly when $\lfloor B/r\rfloor\geq667$, under the stated premises.\\
Proposed federation, $F$, with no additional usable peer service & Same usable service opportunities and non-federation entitlements as $G$ & Same as $G$ & Same as $G$ & No completion gain; positive net overhead makes $F$ more costly on the same accounting basis.\\
\end{longtable}
\par\smallskip{\footnotesize Evidence: [Research session 1, artifact 8]; claims: C2}
\end{landscape}
\textit{Claims discussed: }C2
\section{Evidence limitations and uncertainties}
The evidence limits the central argument to conditional feasibility: it does not establish that federation delivers useful, reliable research access under actual campus policies and budgets. The principal source-grounding gap is consequential rather than merely bibliographic: published capabilities, current API rates, contractual terms and institutional approvals lack inspected supporting evidence. Conditional accounting and authorization arguments therefore cannot establish that their premises hold for a deployable service. [C3; [Research session 1]]

The supplied characterization of MindRouter includes protocol translation, WDRR scheduling, quotas, audit logging and GPU telemetry, but these descriptions are not independently substantiated by the supplied excerpts. The reported review date of 2026-10-31 is not an independently established access date for either the documentation or the PEARC\textquoteright{}26 paper. Federation and frontier routing must consequently remain hypothetical extensions, not attributed product capabilities. The missing source inspection is an evidence limitation, not evidence that the described campus gateway capabilities are absent. [C3]

Missing commercial and institutional evidence also prevents a hard-budget or sensitive-data deployment conclusion. Dated token prices alone would not establish enforceable maximum liability for complete workflows, including retries, tools, embeddings, cache misses and ambiguous charged timeouts; the relevant charging terms remain unestablished. Likewise, destination approval would not by itself establish retention obligations, permission to process derived artifacts, or complete financial mediation. These distinctions separate a proposed admissibility rule from evidence that actual providers and campus agreements satisfy it. [C3]

The financial argument has a specific unresolved recovery weakness, not merely a general need for implementation testing. In the hypothetical stale-backup counterexample, an outstanding liability of $6$ is absent from a restored ledger with a cap of $10$; admitting another liability of $6$ produces total exposure of $12$, even if fresh identifiers prevent replay. Namespace fencing therefore does not alone preserve the budget bound: recovery must also preserve or conservatively cover outstanding exposure. Timeout or cancellation cannot justify releasing a reservation while the original request can still incur charges, and the operational evidence needed to determine final liability remains unspecified. [C3]

The sovereignty argument is similarly narrower than a guarantee against all post-revocation processing. Its stated premise defines revocation at the authorization-decision ordering point, rather than at eventual transmission, receipt or execution. Database fencing cannot by itself stop an already authorized continuation or recall packets beyond campus control, while recipient retention and compliance require separately established obligations. Thus conditional control over an adapter\textquoteright{}s invocation initiation does not establish exactly-once provider execution, post-revocation non-delivery or recipient compliance. The proposed authorization, dispatch and recovery transitions remain unimplemented and unapproved. [C3]

The usefulness calculations expose a separate inferential limitation: aggregate success cannot establish sensitive-task coverage. In the supplied hypothetical cohort, $200q_S+800q_N=800$ is compatible with $q_S=0$ and $q_N=1$, so an aggregate useful-completion rate can coexist with no useful sensitive completions. Assigning sensitive tasks admission priority does not repair that gap. Even an explicit sensitive-task marginal usefulness of $0.80$ would not guarantee realized coverage: the recorded common-shock example has all sensitive tasks succeed with probability $0.80$ and all fail with probability $0.20$. These counterexamples constrain interpretation of the scenario without providing measurements of campus workloads. [C3]

The competing economic interpretations remain conditional for the same reason. A gateway with separately approved external access is cheaper in the recorded case where federation adds no usable service or offsetting benefit but does add positive overhead; this does not establish that peer campuses ordinarily contribute no useful capacity. Conversely, nominal peer GPU capacity does not establish capacity that is authorized, reachable, model-compatible and deadline-feasible for a particular request. Actual campus costs, service capacity, workload performance and net federation benefit remain unknown, so neither general federation superiority nor general federation disadvantage follows. [C3]

The proposed pilot does not close these evidential gaps. Its sampling assumptions and hypothetical planning inputs are not established properties of live campus workloads, and its broader ninety-day protocol and decision thresholds remain incompletely operationalized. A paired comparison would also be limited to tasks authorized for both workflows, rather than automatically representing all requested research access. For reliability and safety, the negative constraints remain important: scheduling weights do not establish unconditional latency guarantees or starvation freedom, and finite fault tests with no observed violations would provide bounded support rather than proof of general safety. [C3]

Taken together, these limitations change the interpretation of the central argument: the record identifies conditions under which controlled routing and budget accounting could be coherent, but does not establish that those conditions are contractually available, institutionally accepted or operationally enforced. The strongest supported comparison is therefore between conditional service designs, not demonstrated campus outcomes. Federation\textquoteright{}s additional complexity has not been justified by measured useful completions, while the simpler gateway-plus-approved-access alternative has not been established as sufficient for actual demand. [C3]

\textit{Claims discussed: }C3
\section{Open questions and next steps}
The next step is to determine whether federation adds a material, policy-permitted research capability beyond a campus gateway with separately approved external access, rather than treating additional remote capacity as sufficient justification. This requires three distinct advances: institutional agreement on permissible processing, an implementation that satisfies the proposed enforcement assumptions, and a workload study demonstrating useful service within actual budgets. None can substitute for the others: an authorization argument does not establish task quality, and greater throughput does not establish affordability or permission to disclose data. [C4]

Source grounding should precede implementation and procurement decisions. The MindRouter documentation and PEARC\textquoteright{}26 paper need to be retrieved, versioned, and recorded with actual access dates so that published gateway capabilities can be distinguished from proposed extensions and independently measured behavior. The project-reported review date of 2026-10-31 is historical context, not evidence of a subsequent inspection; federation, frontier routing, and the proposed financial controls must remain hypothetical unless supported separately. Dated provider prices, charging rules, retention terms, and institutional approvals are also required before using a provider in the comparison. [C4]

The first institutional deliverable should be an approved policy specification covering identity, classification inheritance, destinations, retention, audit, recipient obligations, and cross-campus agreements. It must address derived prompts, outputs, logs, embeddings, and caches, rather than approving only the original input. Security and data owners must also decide whether authorization is effective at the home-campus decision point or whether stronger restrictions on later transmission, receipt, execution, and retention are required; the proposed decision-time semantics permit a previously committed grant to transmit after a later revocation.

Public or synthetic fixtures should therefore precede real-data testing until the relevant campuses approve both the processing domains and these in-flight semantics. [C4]

The implementation task is to establish the assumptions on which the conditional design depends, not merely to reproduce its state labels. Independent inspection should examine complete mediation of egress and spending, protocol canonicalization, immutable request manifests, authoritative ordering of policy and ledger changes, crash-safe one-use issuance, sender fencing, replay-record retention, and durable reconciliation. Translation and attachment resolution must occur before the request is frozen, while status queries, restarts, and duplicate acknowledgements must not reconstruct sending authority.

These are outstanding implementation obligations: no refinement proof or independent inspection currently establishes that the proposed mechanisms satisfy them. [C4]

Fault testing should target the boundaries between authorization, dispatch, and reconciliation, with independent packet and provider observations supplementing protected decision records. The table specifies discriminating tests for the proposed controls; its outcomes are acceptance criteria, not reported results. Financial recovery deserves particular attention because fresh credentials and replay isolation do not restore liabilities omitted by a stale backup, and a current zero-charge response does not establish finality while an original sender can still transmit. Tests should also measure the service lost to safe denial and unresolved reservations, since preserving policy and budget constraints can reduce availability. [C4]

The empirical comparison should retain all four baselines: campus-local inference only, approved frontier API only, a campus gateway with separately approved external access, and that gateway arrangement extended by a hypothetical federation. Each submitted research objective should be registered as one parent task with its input snapshot, rubric, deadline, classification, provenance requirements, and resource envelope; tools, retries, and agent exchanges are descendants rather than additional successes. Overall coverage should use the offered workload, while model-quality comparisons should use matched, policy-eligible strata and report excluded demand explicitly.

Correct refusals belong in the safety record, not in the count of delivered research results. [C4]

The workload needs separate observations for coding, routine analysis, scientific reasoning, multimodal or long-context work, agentic workflows, sensitive processing, and reproducible batches. A useful completion should jointly satisfy authorization, the approved quality rubric, the delivery deadline, and required provenance; recording a model-version label alone is not a reproducibility demonstration. Sensitive-stratum usefulness must be measured directly because aggregate success can conceal complete failure on sensitive tasks, and admission priority does not establish usefulness.

The recorded candidate protocols also propose different deadlines, so a single approved set of task-specific acceptance rules must be frozen before enrollment rather than selected after observing results. [C4]

The federation comparison should hold non-federation entitlements constant, including local allocations, approved paid access, policy rules, and parent-task spending limits, while counting the added partner capacity and agreement costs explicitly. Paired executions should prevent answer sharing and cross-arm cache leakage, with isolated capacity slices or matched replay windows to avoid mistaking interference from duplicate study execution for ordinary-service latency. Queue rejection, model incompatibility, budget exhaustion, outages, and missed deadlines should remain outcomes rather than post hoc exclusions.

Controlled replay and live-service observations should be reported separately because they support different claims about operational behavior. [C4]

Capacity and fairness measurements must identify usable service opportunities, not merely total GPU inventory. Required observations include memory and context fit, compatible capacity, queue blocking, cold starts, maintenance conditions, peak demand, scheduling quanta, and the interaction between batch and interactive work. WDRR weights should not be interpreted as an unconditional latency or starvation-freedom guarantee; the study needs an approved reference allocation and observation window, together with either a demonstrated scheduling guarantee or an explicit report of its absence. Administrative effort and on-call burden should be measured alongside throughput and availability rather than left outside the service comparison. [C4]

The economic study should replace hypothetical ledger entries with campus evidence and preserve the distinction between sunk, incremental, and fully loaded costs. Local accounts need acquisition and replacement cash flows, depreciation, utilization and idle capacity, power and cooling, storage, network, support, procurement, and operations labor, without charging the same replacement or staff activity twice. API accounts need dated rates and contractual bounds covering workflow amplification, retries, tools, embeddings, cache misses, ancillary charges, commitments, and unresolved liabilities; allocation ceilings must not be presented as actual spending.

Sensitivity analysis should vary workload mix, demand, utilization, prices, and federation overhead rather than extrapolate constant completion fractions through saturation. [C4]

A reviewed 90-day pilot requires a decision procedure before enrollment, including a declared primary population, approved rubrics, justified precision, multiplicity handling, fairness and administrative ceilings, and explicit inconclusive outcomes. One proposed strengthening would require the lower endpoint of a prespecified paired confidence interval to exceed $0.10$, rather than accepting a point estimate of at least ten percentage points with an interval that merely excludes zero; this is a candidate decision rule, not established institutional policy. Its sample size and power remain unresolved, and independent-pair calculations cannot be transferred uncritically to correlated live workloads or to a differently weighted block analysis.

Safety stops should end the efficacy claim as inconclusive, while zero observed violations in finite fault tests should remain bounded evidence rather than a general safety proof. [C4]

The resulting decision should ask whether federation delivers an approved capability or useful-completion gain sufficient to justify its additional operating and governance burden, subject to separate safety, coverage, service, and administrative gates. If usable service opportunities are unchanged, there is no offsetting benefit, and federation adds positive overhead on the chosen cost basis, the simpler gateway arrangement is cheaper under those stated conditions. The evidence therefore changes federation from an assumed solution to an extension that must earn deployment through a demonstrated, policy-permitted advantage over campus service plus separately approved external access. [C4]

\begin{landscape}
\small
\begin{longtable}{@{}>{\raggedright\arraybackslash}p{\dimexpr(\linewidth-2\tabcolsep*3)/168*42\relax}>{\raggedright\arraybackslash}p{\dimexpr(\linewidth-2\tabcolsep*3)/168*42\relax}>{\raggedright\arraybackslash}p{\dimexpr(\linewidth-2\tabcolsep*3)/168*42\relax}>{\raggedright\arraybackslash}p{\dimexpr(\linewidth-2\tabcolsep*3)/168*42\relax}@{}}
\caption{Proposed adversarial acceptance tests for the hypothetical enforcement extension. These are prospective criteria under C4, not executed tests or verified MindRouter properties.}\\
\toprule
\textbf{Injected condition} & \textbf{Required observation} & \textbf{Decision consequence} & \textbf{Evidence source}\\
\midrule
\endfirsthead
\toprule
\textbf{Injected condition} & \textbf{Required observation} & \textbf{Decision consequence} & \textbf{Evidence source}\\
\midrule
\endhead
\midrule
\multicolumn{4}{r}{\footnotesize Continued on next page}\\
\endfoot
\bottomrule
\endlastfoot
Partition with identical cached approval, with and without an unobserved home-campus revocation & Observe authorization decisions and packet destinations in paired traces. Under immediate decision-time revocation, fresh disconnected requests must be denied or deferred. & Admission from cached approval fails the stated safety requirement; safe denial must also be counted as lost availability. & [Research session 1, artifact 1]\\
Concurrent revocation or cancellation around dispatch commitment & Establish the authoritative transaction order and the exact attempt authorized. Observe any later transmission separately from the grant decision. & A pre-revocation grant does not establish post-revocation non-delivery; campuses requiring that stronger property need additional controls or must reject these semantics. & [Research session 1, artifact 4]\\
Charged timeout, cancellation acknowledgement, or provisional zero-charge response while delayed sending remains possible & Inspect retained reservations and authoritative evidence of total final liability, including possible delayed requests. & Do not release financial coverage solely because of timeout, cancellation, local silence, or provisional billing evidence. & [Research session 1, artifact 7]\\
Stale-backup restoration with fresh identifiers and credentials & Check both replay isolation and preservation or conservative reconstruction of prior unresolved exposure before new admission. & Fresh namespaces alone are insufficient for financial recovery; missing historical holds must not become spendable capacity. & [Research session 1, artifact 7]\\
Malicious direct egress, unapproved agent-tool destination, ambiguous classification, and simultaneous quota reservations & Independently inspect network destinations, protected authorization records, and all applicable reservation balances under concurrent execution. & Observed bypass, unauthorized disclosure, or reservation overshoot blocks progression to a reviewed live pilot; a finite pass is not deployment authorization. & [Research session 1, artifact 1]\\
\end{longtable}
\par\smallskip{\footnotesize Evidence: [Research session 1, artifact 1], [Research session 1, artifact 4], [Research session 1, artifact 7]; claims: C4}
\end{landscape}
\textit{Claims discussed: }C4
\section{Algorithms, results, and limitations}
\subsection{C1 -- Unresolved}
The supplied domains share authorized root-task accounting and enforceable nested liabilities. The source additionally assumes linearizable home authority, trusted manifest binding, mediated one-use issuance and campus acceptance of decision-time authorization. The historical target original is unavailable. For the repaired expected-count example, each local stratum has average joint authorized quality/deadline/provenance usefulness 0.80, and every selected API cohort has average usefulness 0.90. These are hypothetical probabilities, not observations or guarantees. Every relevant invocation and spending path is mediated by the trusted adapter and ledger, including tools, retries and separately approved accounts. The adapter issues a consumed grant at most once, disables uncertain retransmission and does not recover by reissuing it.
Recorded at portfolio\_decision/cross\_strategy\_reviews/0/transfer\_review/assumption\_comparison; assumptions/2; assumptions/3. This is a recorded research item, not an independently verified conclusion; the cited saved evidence requires assessment.
\textit{Evidence: }[Research session 1], [Research session 1, artifact 8], [Research session 1, artifact 4]
\subsection{C2 -- Unresolved}
Structured argument steps are recorded; review their evidence and objections. 2. Define comparable useful completions before counting costs. Define stratum-specific admissions and usefulness.
Recorded at outcome/argument\_availability/label; investigator\_output/reasoning\_steps/1/title; investigator\_output/reasoning\_steps/1/title. This is a recorded research item, not an independently verified conclusion; the cited saved evidence requires assessment.
\textit{Evidence: }[Research session 1], [Research session 1, artifact 2], [Research session 1, artifact 8]
\subsection{C3 -- Limitation}
Published capabilities, current rates, contractual terms and campus approvals lack inspected source evidence.
Recorded at outcome/caveats/2. This is a recorded research item, not an independently verified conclusion; the cited saved evidence requires assessment.
\textit{Evidence: }[Research session 1]
\subsection{C4 -- Unresolved}
o96/o102/o125/o130 and earlier comparison duties remain open: complete the four-baseline workload, quality, provenance, latency, throughput, GPU-fit, fairness, administrative-burden and cost study, with dated prices, sunk/incremental/fully loaded accounting, budget sensitivities and a reviewed 90-day pilot. o38/o45/o50/o54 remain unresolved: approve classification inheritance, destinations, identity, retention, audit, revocation and in-flight semantics, recipient obligations and cross-campus agreements; implement and test the proposed transitions before making safety claims. The specification advances o3/o29/o32 but does not discharge implementation obligations: verify complete mediation, protocol canonicalization, immutable manifests, linearizability, crash-safe one-use issuance, sender fencing, replay-record lifetime and durable reconciliation. No implementation refinement proof or independent inspection exists.
Recorded at investigator\_output/open\_obligations/4; investigator\_output/open\_obligations/4; investigator\_output/open\_obligations/2. This is a recorded research item, not an independently verified conclusion; the cited saved evidence requires assessment.
\textit{Evidence: }[Research session 1, artifact 7], [Research session 1, artifact 5], [Research session 1, artifact 4]
\section{Limitations}
\begin{itemize}
\item The saved record may contain unverified interpretations and unresolved evidence conflicts; these require independent assessment.
\end{itemize}
\section{Open obligations}
\begin{itemize}
\item Retrieve both starting sources, record actual access dates and versions, and verify the user-supplied characterization of protocol translation, WDRR, quotas, audit logging, telemetry, and on-campus scope. No federation or frontier-routing feature has been verified.
\item Retrieve the MindRouter documentation and PEARC'26 paper, record actual access dates and versions, and verify the supplied feature descriptions and any changes. Obtain dated API rates, contractual charging bounds, retention terms and institutional approval; none is established here.
\item Inspect both starting references, record actual access dates and versions, and verify the precise published MindRouter capabilities and any changed claims. Do not infer proposed federation or financial controls from the supplied gateway description.
\item Source grounding remains unresolved: obtain both starting references, record actual access dates and versions, and verify protocol translation, WDRR, quotas, audit logging, telemetry and campus-only scope. The project-reported 2026-10-31 review date is not this session's access date; no changed claim has been verified.
\item o34/o41/o51/o55 remain unresolved: retrieve both starting sources and their original saved context where needed, record actual dates/versions, verify published features and changes, and obtain dated provider prices, contracts, retention and approvals. Peer extracts are truncated and their omitted details remain unknown.
\item Verify the MindRouter documentation and PEARC'26 paper with actual access dates and versions; inspect saved originals where needed. Source obligations o7/o56/o63/o68 remain open.
\item o94/o100/o122/o127 and related implementation duties remain open: implement and independently inspect the transaction schema, immutable translation, sender continuation, retry suppression, replay retention, restore fencing and complete mediation. The present sketch advances specification only.
\item Retrieve the two starting references and pointer-only originals before reusing their detailed claims or numerical ledger; record actual access dates/versions and obtain dated API prices, charging bounds, retention terms and institutional approvals. Source-related ledger obligations remain unresolved.
\item Ask campus security and data owners to define classification inheritance, approved domains, revocation timing, in-flight handling, and acceptable lease semantics; ask procurement/legal to validate partner/provider retention, subprocessors, audit rights, and enforceable charge bounds.
\item Collect a campus cost ledger that separates sunk, incremental and fully loaded cost, including replacement timing, power/cooling, storage/network, support, procurement and staff effort. Replace illustrative aggregate totals before making a procurement recommendation.
\item Obtain campus security/legal decisions on data classification, derived artifacts, identity federation, retention, incident response, contracts and permitted domains. Determine whether remote execution can actually enforce them.
\item Campus/security/legal/procurement decisions remain required for classification inheritance, decision-time versus stronger revocation semantics, in-flight handling, destination and retention approval, final billing evidence and enforceable liability bounds.
\item o35/o42 remain unresolved: replace every hypothetical ledger row with campus evidence, validate avoidability and staff scopes, inventory heterogeneous GPUs, and reconcile purchase/replacement cash flow with depreciation and stranded capacity.
\item Obtain campus security/legal/procurement approvals, enforceable billing bounds and complete financial/egress mediation; test policy freshness, retention, cancellation and recovery. Obligations o8\textendash{}o10/o57\textendash{}o58/o65\textendash{}o69 remain open.
\item o95/o101 and related testing duties remain open: execute the fault schedules and independent egress/provider-observation tests, including adversarial tenants and stale credentials; quantify availability and fairness losses from uncertainty holds.
\item Replace hypothetical costs with campus evidence, preserving sunk, incremental, fully loaded and acquisition/replacement cash-flow distinctions. Validate federation overhead and any offsetting benefit rather than assuming them.
\end{itemize}
\section{Frozen obligation reconciliation}
Unlisted obligations are not presumed resolved. The archived research snapshot retains complete source references and the full reconciliation.
\begin{itemize}
\item Retrieve the MindRouter documentation and PEARC'26 paper, record actual access dates and versions, and verify the supplied feature descriptions and any changes. Obtain dated API rates, contractual charging bounds, retention terms and institutional approval; none is established here. \textit{ (Listed as open in paper; 7 saved source reference(s))}
\item Collect a campus cost ledger that separates sunk, incremental and fully loaded cost, including replacement timing, power/cooling, storage/network, support, procurement and staff effort. Replace illustrative aggregate totals before making a procurement recommendation. \textit{ (Listed as open in paper; 7 saved source reference(s))}
\item Measure matched task quality, complete workflow token/call amplification, latency distributions, fit/cold-start behavior, capacity, unserved demand and provenance limits. Estimate uncertainty and sensitivity to task mix, utilization, API prices and demand rather than extrapolating constant success fractions through saturation. \textit{ (Not listed by author; not presumed resolved; 7 saved source reference(s))}
\item Establish whether all paid calls, tools and retries can be assigned enforceable maximum charges and whether timeout reconciliation and nested quotas preserve the reservation invariant under concurrency and failures. This requires implementation testing and contract review. \textit{ (Not listed by author; not presumed resolved; 7 saved source reference(s))}
\item Resolve campus identity, approved-domain classification, derived-artifact handling, retention, audit, incident and cross-campus agreement requirements before any sensitive or federated testing. \textit{ (Not listed by author; not presumed resolved; 7 saved source reference(s))}
\item Run the bounded pilot and test whether federation adds enough useful completions to justify its net overhead. Safety fault injection with zero observed violations is only bounded evidence, not proof of general safety. \textit{ (Not listed by author; not presumed resolved; 7 saved source reference(s))}
\item Ask campus computing about measured capacity and on-call effort; security about classification and derived-data boundaries; procurement about price changes, commitments and ambiguous charges; faculty about quality and reproducibility; students about access, deadlines and fairness. Their answers may reverse the preferred baseline. \textit{ (Not listed by author; not presumed resolved; 7 saved source reference(s))}
\item o15: Retrieve and date the starting sources; obtain current rate, charging, retention and approval evidence. \textit{ (Not listed by author; not presumed resolved; 8 saved source reference(s))}
\item o16: Replace aggregate illustrative costs with a non-overlapping campus ledger separating sunk, incremental and fully loaded costs. \textit{ (Not listed by author; not presumed resolved; 8 saved source reference(s))}
\item o17: Establish matched quality, workflow amplification, latency, capacity, provenance and sensitivity evidence. \textit{ (Not listed by author; not presumed resolved; 8 saved source reference(s))}
\item o18: Establish enforceable charge bounds and test concurrent reservations, reconciliation and nested quotas. \textit{ (Not listed by author; not presumed resolved; 8 saved source reference(s))}
\item o19: Resolve classification, derived-artifact boundaries, identity, retention, audit and cross-campus authority. \textit{ (Not listed by author; not presumed resolved; 8 saved source reference(s))}
\item o20: Execute the approved bounded pilot and assess federation's net benefit with uncertainty and independent safety gates. \textit{ (Not listed by author; not presumed resolved; 8 saved source reference(s))}
\item o21: Obtain computing, security, procurement, faculty and student input. \textit{ (Not listed by author; not presumed resolved; 8 saved source reference(s))}
\item Material, recoverable scope gap: an accounting threshold and aggregate synthetic example do not yet provide the requested four-baseline cost, quality, coverage and capacity comparison. \textit{ (Not listed by author; not presumed resolved; 6 saved source reference(s))}
\item Material, recoverable implementation gap: the budget invariant is conditional on enforceable liability bounds and atomic nested-ledger transitions. It is not evidence that provider billing, cancellations, timeouts or concurrent quotas satisfy those premises. \textit{ (Not listed by author; not presumed resolved; 6 saved source reference(s))}
\item Material, recoverable safety gap: the architecture does not yet specify policy freshness/revocation and reservation ownership across admission, dispatch, retries and reconciliation. No deployment-safety conclusion follows. \textit{ (Not listed by author; not presumed resolved; 6 saved source reference(s))}
\item Evidence limitation, not a contradiction: source descriptions, current rates, retention terms, campus policy and workload performance remain unverified. \textit{ (Not listed by author; not presumed resolved; 6 saved source reference(s))}
\end{itemize}
213 further distinct obligations are retained in the archived reconciliation.
\section{Conclusions}
The indexed record supports only provisional interpretation: C1 (The supplied domains share authorized root-task accounting and enforceable nested liabilities. The source additionally assumes linearizable home authority, trusted manifest binding, mediated one-use issuance and campus acceptance of decision-time authorization. The historical target original is unavailable. For the repaired expected-count example, each local stratum has average joint authorized quality/deadline/provenance usefulness 0.80, and every selected API cohort has average usefulness 0.90. These are hypothetical probabilities, not observations or guarantees.

Every relevant invocation and spending path is mediated by the trusted adapter and ledger, including tools, retries and separately approved accounts. The adapter issues a consumed grant at most once, disables uncertain retransmission and does not recover by reissuing it.); C2 (Structured argument steps are recorded; review their evidence and objections. 2. Define comparable useful completions before counting costs. Define stratum-specific admissions and usefulness.); C3 (Published capabilities, current rates, contractual terms and campus approvals lack inspected source evidence.). Unresolved evidence and obligations remain; no independent verification is claimed.

\textit{Claims discussed: }C1, C2, C3
\appendix
\section{Selected source-linked research excerpts}
The following are bounded, selected excerpts of original research records. They are not exhaustive; consult the full frozen snapshot archived with this paper. They do not constitute independent verification.
\subsection{Research excerpt 1}
\textit{Source and field: }[Research session 1, artifact 3]: skeptic\_output/reasoning\_assessments\par
[\{"basis\_artifact\_ids": [], "disposition": "unresolved", "objection": "The entry is labeled an assumption, but the published-capability premise is not one of the five explicit conditions of the ledger claim. Retain it as unverified background, not an established premise. The ledger proof does not require it.", "rationale": "The attribution matches the supplied project description, and the separation of published gateway claims from proposed extensions is appropriate.

No original source establishes those capabilities here.", "step\_number": 1\}, \{"basis\_artifact\_ids": [], "disposition": "accepted", "objection": "", "rationale": "As a proposed admission rule, constructing eligibility before ranking prevents a better score or failure from authorizing a prohibited destination. Per-call re-evaluation and mediation of derived artifacts address the stated policy scope.", "step\_number": 2\}, \{"basis\_artifact\_ids": [], "disposition": "accepted", "objection": "", "rationale": "Atomic admission against every applicable scope preserves each inequality.

Attributing the same liability to nested scopes does not imply multiple billing, and separating GPU quotas from money avoids conflating resources.", "step\_number": 3\}, \{"basis\_artifact\_ids": [], "disposition": "accepted", "objection": "", "rationale": "The settlement change is z-r \textless{}= 0; unresolved events leave totals unchanged. Together with the initial invariant and the listed durability and accounting-window assumptions, this proves allocation safety. The explicit cash-flow limitation prevents overclaiming invoice-month control.",\ldots{}

\subsection{Research excerpt 2}
\textit{Source and field: }[Research session 1]: steps/4/output/reasoning\_assessments\par
[\{"basis\_artifact\_ids": [], "disposition": "unresolved", "objection": "The entry is labeled an assumption, but the published-capability premise is not one of the five explicit conditions of the ledger claim. Retain it as unverified background, not an established premise. The ledger proof does not require it.", "rationale": "The attribution matches the supplied project description, and the separation of published gateway claims from proposed extensions is appropriate.

No original source establishes those capabilities here.", "step\_number": 1\}, \{"basis\_artifact\_ids": [], "disposition": "accepted", "objection": "", "rationale": "As a proposed admission rule, constructing eligibility before ranking prevents a better score or failure from authorizing a prohibited destination. Per-call re-evaluation and mediation of derived artifacts address the stated policy scope.", "step\_number": 2\}, \{"basis\_artifact\_ids": [], "disposition": "accepted", "objection": "", "rationale": "Atomic admission against every applicable scope preserves each inequality.

Attributing the same liability to nested scopes does not imply multiple billing, and separating GPU quotas from money avoids conflating resources.", "step\_number": 3\}, \{"basis\_artifact\_ids": [], "disposition": "accepted", "objection": "", "rationale": "The settlement change is z-r \textless{}= 0; unresolved events leave totals unchanged. Together with the initial invariant and the listed durability and accounting-window assumptions, this proves allocation safety. The explicit cash-flow limitation prevents overclaiming invoice-month control.",\ldots{}

\subsection{Research excerpt 3}
\textit{Source and field: }[Research session 1, artifact 1]: skeptic\_output/reasoning\_assessments\par
[\{"basis\_artifact\_ids": [], "disposition": "accepted", "objection": "", "rationale": "Accepted as the exact conditional model, not as established university policy. It explicitly excludes alternative information channels and previously committed request-specific authorization, and defines safety and availability sufficiently for the argument.", "step\_number": 1\}, \{"basis\_artifact\_ids": [], "disposition": "accepted", "objection": "", "rationale": "An unobserved home-campus revocation can be inserted without changing the admission point's request, credentials, clock or cached state.

A signature on historical policy provides authenticity, not knowledge that no later revocation occurred.", "step\_number": 2\}, \{"basis\_artifact\_ids": [], "disposition": "accepted", "objection": "", "rationale": "Identical observations imply identical deterministic decisions. Coupling random choices gives equal admission probability in the two executions, so zero unauthorized-admission probability excludes positive partition availability under these assumptions.", "step\_number": 3\}, \{"basis\_artifact\_ids": [], "disposition": "accepted", "objection": "", "rationale": "Campus-approved non-revocable leases weaken the immediate-revocation requirement, while an authoritative connected decision service removes the disconnected-authorizer premise.

The argument correctly does not establish recall or prevention of delivery after an earlier authorization.", "step\_number": 4\}, \{"basis\_artifact\_ids": [], "disposition": "accepted", "objection": "", "rationale": "Accepted as a proposed conservative policy\ldots{}

\subsection{Research excerpt 4}
\textit{Source and field: }[Research session 1, artifact 4]: skeptic\_output/reasoning\_assessments\par
[\{"basis\_artifact\_ids": ["[Research session 1, artifact 1]", "[Research session 1, artifact 2]", "[Research session 1, artifact 3]"], "disposition": "accepted", "objection": "", "rationale": "The supplied prior extract identifies the stated policy and reservation gaps. Both peer extracts explicitly list liability-bound and complete-mediation premises. The step correctly treats these as unverified inputs and does not import the truncated economic argument.", "step\_number": 1\}, \{"basis\_artifact\_ids": [], "disposition": "accepted", "objection": "", "rationale": "Accepted as conditional model definitions.

Separating root task, immutable logical operation and potentially billable attempt prevents retries from becoming extra useful completions and prevents authorization reuse after changing a destination or payload.", "step\_number": 2\}, \{"basis\_artifact\_ids": [], "disposition": "accepted", "objection": "", "rationale": "Accepted as an explicit authority and campus-acceptance premise, not established policy.

Linearization gives revocation and grant a definite order while expressly allowing later transmission from an earlier grant; disconnected cached approval cannot create a new commit.", "step\_number": 3\}, \{"basis\_artifact\_ids": [], "disposition": "accepted", "objection": "", "rationale": "Accepted as the conditional state and trusted-validation specification. Separating financial state from task outcome avoids inferring settlement from success, failure or cancellation. Unknown classifications and unsupported request semantics fail closed rather than being silently\ldots{}

\subsection{Research excerpt 5}
\textit{Source and field: }[Research session 1, artifact 6]: skeptic\_output/reasoning\_assessments\par
[\{"basis\_artifact\_ids": ["[Research session 1, artifact 3]"], "disposition": "accepted", "objection": "", "rationale": "The supplied cycle-1 extract contains the stated reservation and scheduling examples. The current proposal explicitly tightens, rather than retrospectively reinterprets, the earlier federation gate and preserves its open requirements.", "step\_number": 1\}, \{"basis\_artifact\_ids": ["[Research session 1, artifact 3]"], "disposition": "accepted", "objection": "", "rationale": "As a proposed comparison, matching non-federation entitlements isolates the addition of remote capacity.

Explicit financial mediation and policy-first eligibility avoid treating approval or failure as sufficient authorization.", "step\_number": 2\}, \{"basis\_artifact\_ids": [], "disposition": "accepted", "objection": "", "rationale": "A common preregistered paired denominator prevents admission and completion filtering from inflating arm success. Reporting requested, authorized, paired and admitted counts exposes eligibility selection and limits the estimand's scope.", "step\_number": 3\}, \{"basis\_artifact\_ids": [], "disposition": "accepted", "objection": "", "rationale": "Success is conjunctive in quality, provenance, authorization and delivery deadline. End-to-end latency includes workflow delays rather than only generation time.

The numerical deadlines and incomplete pass rules are correctly presented as proposals awaiting approval.", "step\_number": 4\}, \{"basis\_artifact\_ids": [], "disposition": "accepted", "objection": "", "rationale": "The pairing, cache isolation and matched-load\ldots{}

\subsection{Research excerpt 6}
\textit{Source and field: }[Research session 1, artifact 7]: skeptic\_output/reasoning\_assessments\par
[\{"basis\_artifact\_ids": [], "disposition": "unresolved", "objection": "The claim that supplied source text provides the earlier states cannot be checked here. Request the relevant definitions only; do not treat a pointer or ledger summary as the original argument.", "rationale": "The evidence boundary and retained decision-time and timeout-release limitations are appropriate. However, this review packet contains pointer-only prior artifacts and ledger summaries rather than the described cycle-2 state-machine extract.", "step\_number": 1\}, \{"basis\_artifact\_ids": [], "disposition": "accepted", "objection": "", "rationale": "Accepted as a proposed specification premise, not an implemented property.

The records bind identity, destination and frozen content, and the explicit classification premise is necessary because a digest alone cannot establish policy correctness.", "step\_number": 2\}, \{"basis\_artifact\_ids": [], "disposition": "accepted", "objection": "", "rationale": "Accepted as a conditional schema and transaction-domain premise. Unique reservation and settlement keys plus exact accounting units support deterministic admission accounting, provided the stated serialization and complete-scope coverage actually hold.", "step\_number": 3\}, \{"basis\_artifact\_ids": [], "disposition": "accepted", "objection": "", "rationale": "Under the stated atomic transaction premise, commit and pre-send closure cannot both win the same state predicate.

Releasing holds only in the winning closure transaction prevents a partially released committed reservation.", "step\_number": 4\},\ldots{}

\subsection{Research excerpt 7}
\textit{Source and field: }[Research session 1, artifact 2]: skeptic\_output/reasoning\_assessments\par
[\{"basis\_artifact\_ids": [], "disposition": "accepted", "objection": "", "rationale": "The feature attribution matches the supplied project description, and the four baseline definitions preserve the distinction between the published campus gateway and hypothetical extensions. The stated 2026-10-31 date is correctly identified as supplied history, not this session's access date.", "step\_number": 1\}, \{"basis\_artifact\_ids": [], "disposition": "accepted", "objection": "", "rationale": "The binary root-task definition entails S\_j as a count of authorized, timely, rubric-passing completions. Keeping refusals and ineligible demand visible prevents treating policy exclusion as successful service.

Separating quality strata and provenance avoids equating call count or version recording with useful, reproducible work.", "step\_number": 2\}, \{"basis\_artifact\_ids": [], "disposition": "accepted", "objection": "", "rationale": "Accepted as a proposed experimental specification under the listed assumption that pilot deadlines, samples and gates require campus approval. The task-specific rubrics and latency decomposition are suitable candidate controls; no achieved service level or model equivalence is inferred.", "step\_number": 3\}, \{"basis\_artifact\_ids": [], "disposition": "accepted", "objection": "", "rationale": "The annualized ownership formula is dimensionally consistent when power is in kW and electricity is priced per kWh.

Cooling and replacement double-counting are explicitly excluded. Distinguishing consolidated transfers, campus budgets, sunk costs and distinct active-researcher\ldots{}

\subsection{Research excerpt 8}
\textit{Source and field: }[Research session 1]: outcome/solution\_review/strategy\_findings/0/reasoning\_steps\par
[\{"argument": "The supplied extract of [Research session 1, artifact 4] provides the cycle-2 states and expected counterexamples; its original is not supplied and parts of the extract are incomplete. I use that extract as an unverified specification to refine, not a certified premise. The original [Research session 1, artifact 1] is pointer-only and has not been inspected. Retain two explicit constraints from the supplied cycle-2 text: a pre-revocation grant can transmit later under decision-time semantics; releasing on timeout can expose two liabilities of 6 against cap 10.

The present design neither restores the stronger delivery guarantee nor reuses timeout release.", "basis\_artifact\_ids": ["[Research session 1, artifact 4]"], "epistemic\_status": "deduction", "number": 1, "research\_artifact\_id": "[Research session 1, artifact 7]", "review": \{"basis\_artifact\_ids": [], "disposition": "unresolved", "objection": "The claim that supplied source text provides the earlier states cannot be checked here. Request the relevant definitions only; do not treat a pointer or ledger summary as the original argument.", "rationale": "The evidence boundary and retained decision-time and timeout-release limitations are appropriate.

However, this review packet contains pointer-only prior artifacts and ledger summaries rather than the described cycle-2 state-machine extract."\}, "title": "1. Evidence boundary and retained negative constraints"\}, \{"argument": "Proposed protected tables: Operation(home\_id, issuer\_epoch, root\_task\_id, operation\_id, principal, project, department,\ldots{}

\subsection{Research excerpt 9}
\textit{Source and field: }[Research session 1, artifact 6]: objective\_review/proof\_obligations\par
["Conditional reservation invariant: supported by cycle 1, artifact [Research session 1, artifact 3], supplied reasoning steps 3\textendash{}4; implementation and contractual instantiation remain open.", "Negative constraints on timeout release and unconditional scheduling guarantees: supported by cycle 1, artifact [Research session 1, artifact 3], supplied reasoning steps 5 and 7.", "Paired estimand and variance: supported analytically by cycle 2, current reasoning step 6.", "Conservative interval coverage and multiplicity control: supported under the stated sampling models by cycle 2, current reasoning steps 7 and 11.", "Illustrative precision limits and independent-block bound: supported by cycle 2, current reasoning steps 8\textendash{}10; these are hand-derived planning results, not executed computations.", "Actual gate-power guarantee, fixed affordable enrollment and defensible sampling units: open in cycle 2.", "Operational approvals, sourced costs, implementation safety and measured useful access: open across cycles 1\textendash{}2."]

\subsection{Research excerpt 10}
\textit{Source and field: }[Research session 1, artifact 7]: objective\_review/proof\_obligations\par
["Cycle 3 steps 4\textendash{}7 support the conditional uniqueness argument; implementation must establish their transaction, acknowledgement and one-shot-continuation premises.", "Cycle 3 steps 10, 12\textendash{}14 support conditional coverage during ambiguity; an operational finality predicate must cover delayed sending and all applicable financial scopes.", "Cycle 3 step 9 must be corrected so recovery preserves both replay rejection and unresolved financial exposure.", "Exact refinement correspondence requires the absent cycle-2 definitions; no prior-cycle proof obligation is discharged from its pointer.", "Project achievement additionally requires the source-grounded four-baseline comparison, governance review, transparent dated costs, measurements and bounded pilot retained in the ledger."]

\subsection{Research excerpt 11}
\textit{Source and field: }[Research session 1, artifact 4]: skeptic\_output/obligation\_reviews/0/basis\_artifact\_ids\par
["[Research session 1, artifact 1]"]

9 further research records are retained in the archived source snapshot.
\section*{Evidence and authorship note}
Snapshot SHA-256: 552f9458e600d3e7d0b3aa6c1dd105cc62dab4967db652123d2f134783b76f34
Generated by an AI author from the cited saved research snapshot. This document is not independent verification. No external bibliography was added.
\section*{Validation warnings}
\begin{itemize}
\item Evidence indexing omitted 9706 material saved research field fragment(s). Those fragments were not reviewed for this paper; the interpretation is provisional where the omitted evidence may matter. Consult the frozen source snapshot and omission ledger for full details.
\item The AI planning stage did not yield a usable outline. A conservative, source-linked outline was reconstructed from saved evidence without repeating the planning request; review the resulting paper's structure.
\item The evidence index explicitly omitted 9706 saved field fragment(s); see the generation coverage ledger and review the complete archived snapshot.
\item Structured paper indexing selected 390 source-linked text fragments from 13 archived sources without AI extraction calls. The complete frozen snapshot remains archived; this bounded index does not constitute review of every saved field.
\item Citation-like external references appear in the draft, but this paper has no verified external-source catalog; treat them as unsupported.
\item The frozen evidence contains open proof/research obligation(s) requiring reconciliation: 215 item(s) are not listed by the author and are not presumed resolved; 0 author-listed item(s) have no exact frozen-record match. See the obligation reconciliation.
\item The methods narrative has no linked claim IDs; check it for unsupported factual statements.
\item AI-authored analysis is not independent formal verification; model agreement does not establish proof.
\end{itemize}
\section*{Typesetting verification note}
Needs review: TeX verification is incomplete. The authored research content and safe accepted typesetting corrections were retained. Review the saved proofreading diagnostics before relying on the layout.
\end{document}
